\documentclass[11pt,a4paper]{article}
% ══════════════════════════════════════════════════════════════════════════════
% Edexcel Further Mathematics (Core Pure 1) --- Shared LaTeX Preamble
% Elegant, Minimalist & Aesthetic Academic Design System
% ══════════════════════════════════════════════════════════════════════════════

\usepackage[a4paper, top=22mm, bottom=22mm, left=20mm, right=20mm]{geometry}
\usepackage{amsmath, amssymb, amsthm}
\usepackage{enumitem}
\usepackage{booktabs}
\usepackage{array}
\usepackage{parskip}
\usepackage{microtype}
\usepackage{fancyhdr}
\usepackage{titlesec}
\usepackage{mdframed}
\usepackage{xcolor}
\usepackage[hidelinks]{hyperref}
\usepackage{graphicx}
\usepackage{tikz}
\usepackage{pgfplots}
\pgfplotsset{compat=1.18}

% ── Refined Palette ───────────────────────────────────────────────────────────
\definecolor{primary}{gray}{0.10}
\definecolor{darkgray}{gray}{0.20}
\definecolor{medgray}{gray}{0.50}
\definecolor{lightgray}{gray}{0.96}
\definecolor{boxgray}{gray}{0.85}
\definecolor{accentline}{gray}{0.75}

% ── Section formatting ──────────────────────────────────────────────────────────
\titleformat{\section}[block]
  {\normalfont\bfseries\color{primary}}{}
  {0pt}{}[\vspace{1pt}{\color{accentline}\rule{\textwidth}{0.4pt}}\vspace{4pt}]
\titlespacing*{\section}{0pt}{14pt}{6pt}

% ── List formatting ─────────────────────────────────────────────────────────────
\setlist[enumerate,1]{label=\textbf{\arabic*.}, leftmargin=1.8em, itemsep=14pt}
\setlist[enumerate,2]{label=\textbf{(\alph*)}, leftmargin=1.8em, itemsep=6pt}

% ── Header / footer (Overlaps permanently eliminated) ─────────────────────────
\pagestyle{fancy}
\newcommand{\CPHeader}[2]{%
  \fancyhf{}%
  \setlength{\headheight}{22pt}%
  \renewcommand{\headrulewidth}{0.3pt}%
  \renewcommand{\headrule}{\hbox to\headwidth{\color{accentline}\leaders\hrule height \headrulewidth\hfill}}%
  \fancyhead[L]{\footnotesize\scshape\color{medgray} Core Pure 1: #1}%
  \fancyhead[R]{\footnotesize\color{medgray} Chapter \##2}%
  \renewcommand{\footrulewidth}{0.3pt}%
  \renewcommand{\footrule}{\hbox to\headwidth{\color{accentline}\leaders\hrule height \footrulewidth\hfill}}%
  \fancyfoot[L]{\footnotesize\color{medgray} Pearson Edexcel A-Level Further Mathematics $\cdot$ Core Pure 1}%
  \fancyfoot[C]{}% No center text to ensure zero collision with long titles
  \fancyfoot[R]{\footnotesize\color{medgray} \textbf{Henry.K} \enspace$\vert$\enspace Page \thepage}%
}

% ── Elegant Aesthetic Title Block ──────────────────────────────────────────────
\newcommand{\CPTitleBlock}[2]{%
  \vspace*{-8pt}
  \begin{center}
    {\footnotesize\scshape\color{medgray} Pearson Edexcel A-Level Further Mathematics \enspace$\cdot$\enspace Core Pure 1}\\[4pt]
    {\LARGE\bfseries\color{primary} #1}\\[4pt]
    {\small\color{darkgray} #2}\\[6pt]
    {\footnotesize\color{medgray} Problem Series authored by \textbf{Henry.K} \quad$\cdot$\quad 10 Exam Archetypes \quad$\cdot$\quad 65 Total Marks}\\[8pt]
    {\color{accentline}\rule{\textwidth}{0.5pt}}
  \end{center}
  \vspace{-6pt}
  \noindent{\footnotesize\itshape \textbf{Instructions:} Answer all ten questions. Full algebraic working must be shown for method marks; unsupported answers from calculators will not gain full credit. Give exact values in surd or fractional form unless specified otherwise.}\par
  \vspace{3pt}
  {\color{accentline}\rule{\textwidth}{0.5pt}}
  \vspace{8pt}
}

% ── Answer / Method / Investigate-further boxes ────────────────────────────────
\newmdenv[
  backgroundcolor=lightgray,
  linecolor=boxgray,
  linewidth=0.5pt,
  roundcorner=3pt,
  leftmargin=0pt, rightmargin=0pt,
  innerleftmargin=10pt, innerrightmargin=10pt,
  innertopmargin=6pt, innerbottommargin=6pt,
  skipabove=5pt, skipbelow=5pt
]{ansbox}

\newmdenv[
  backgroundcolor=lightgray!40,
  linecolor=darkgray,
  linewidth=1.2pt,
  leftline=true, rightline=false, topline=false, bottomline=false,
  leftmargin=0pt, rightmargin=0pt,
  innerleftmargin=10pt, innerrightmargin=8pt,
  innertopmargin=5pt, innerbottommargin=5pt,
  skipabove=5pt, skipbelow=7pt
]{invbox}

\newcommand{\ans}[1]{%
  \begin{ansbox}
    \textbf{Answer:} #1
  \end{ansbox}}

\newcommand{\method}[1]{%
  {\small\textbf{Method Criteria:}\par\vspace{2pt}#1}\par}

\newcommand{\inv}[1]{%
  \begin{invbox}
    \textbf{\small Deep Investigation:} {\small #1}
  \end{invbox}}

\newcommand{\qn}[1]{\textbf{#1.}\;}

\CPHeader{Proof by Mathematical Induction --- Solutions}{8}

\begin{document}
\thispagestyle{plain}

\CPTitleBlock{Proof by Mathematical Induction}{Solutions \& Mark Scheme \#8 \enspace$\cdot$\enspace Summations, Divisibility, Matrix Powers \& Recurrence Relations}

\section*{Section A \quad Procedural \& Direct Foundations}

%── Q1 ────────────────────────────────────────────────────────────────
\qn{1} \textbf{Summation of $r(r+1)$}

\ans{Proof complete. $\displaystyle\sum_{r=1}^{n} r(r+1) = \dfrac{n(n+1)(n+2)}{3}$ for all $n \geq 1$.}

\method{
\textbf{Step 1 — Basis case ($n = 1$):}
\begin{itemize}[leftmargin=1.5em,itemsep=2pt]
  \item \textbf{B1}: LHS $= 1 \cdot 2 = 2$. \quad RHS $= \dfrac{1 \cdot 2 \cdot 3}{3} = 2$. \quad LHS $=$ RHS. \checkmark
\end{itemize}

\textbf{Step 2 — Assumption ($n = k$):}
\begin{itemize}[leftmargin=1.5em,itemsep=2pt]
  \item \textbf{M1}: Assume $\displaystyle\sum_{r=1}^{k} r(r+1) = \dfrac{k(k+1)(k+2)}{3}$ holds for some positive integer $k$.
\end{itemize}

\textbf{Step 3 — Inductive step ($n = k+1$):}
\begin{itemize}[leftmargin=1.5em,itemsep=2pt]
  \item \textbf{M1}: Consider $\displaystyle\sum_{r=1}^{k+1} r(r+1) = \dfrac{k(k+1)(k+2)}{3} + (k+1)(k+2)$.
  \item \textbf{M1}: Factor out $(k+1)(k+2)$:
  \[
  = (k+1)(k+2)\!\left(\frac{k}{3} + 1\right) = (k+1)(k+2)\cdot\frac{k+3}{3} = \frac{(k+1)(k+2)(k+3)}{3}\,.
  \]
  \item \textbf{A1}: This is the formula with $n = k+1$. \checkmark
\end{itemize}

\textbf{Step 4 — Conclusion:}
\begin{itemize}[leftmargin=1.5em,itemsep=2pt]
  \item \textbf{B1}: True for $n=1$; true for $n = k+1$ whenever true for $n = k$. By the principle of mathematical induction, true for all $n \geq 1$.
\end{itemize}
}

\inv{The result $\sum_{r=1}^n r(r+1) = \frac{n(n+1)(n+2)}{3}$ counts something combinatorial: specifically, $3\binom{n+2}{3}$, the number of ways to choose 3 items from $n+2$ with an ordering. Can you express $\sum r(r+1)(r+2)$ in similar form? (Hint: it equals $\frac{n(n+1)(n+2)(n+3)}{4}$.)}

\vspace{10pt}

%── Q2 ────────────────────────────────────────────────────────────────
\qn{2} \textbf{Sum of $r \cdot 2^r$}

\ans{Proof complete. $\displaystyle\sum_{r=1}^{n} r \cdot 2^r = 2 + (n-1)\cdot 2^{n+1}$ for all $n \geq 1$.}

\method{
\textbf{Step 1 — Basis case ($n = 1$):}
\begin{itemize}[leftmargin=1.5em,itemsep=2pt]
  \item \textbf{B1}: LHS $= 1 \cdot 2^1 = 2$. \quad RHS $= 2 + (0)\cdot 2^{2} = 2$. \quad LHS $=$ RHS. \checkmark
\end{itemize}

\textbf{Step 2 — Assumption ($n = k$):}
\begin{itemize}[leftmargin=1.5em,itemsep=2pt]
  \item \textbf{M1}: Assume $\displaystyle\sum_{r=1}^{k} r \cdot 2^r = 2 + (k-1)\cdot 2^{k+1}$ for some $k \geq 1$.
\end{itemize}

\textbf{Step 3 — Inductive step ($n = k+1$):}
\begin{itemize}[leftmargin=1.5em,itemsep=2pt]
  \item \textbf{M1}: $\displaystyle\sum_{r=1}^{k+1} r \cdot 2^r = 2 + (k-1)\cdot 2^{k+1} + (k+1)\cdot 2^{k+1}$.
  \item \textbf{M1}: Collect the $2^{k+1}$ terms:
  \[
  = 2 + 2^{k+1}\!\bigl[(k-1)+(k+1)\bigr] = 2 + 2^{k+1} \cdot 2k = 2 + k \cdot 2^{k+2}\,.
  \]
  \item \textbf{A1}: Required form for $n = k+1$: $2 + \bigl((k+1)-1\bigr)\cdot 2^{(k+1)+1} = 2 + k \cdot 2^{k+2}$. \checkmark
\end{itemize}

\textbf{Step 4 — Conclusion:}
\begin{itemize}[leftmargin=1.5em,itemsep=2pt]
  \item \textbf{B1}: By mathematical induction, the result holds for all $n \geq 1$.
\end{itemize}
}

\inv{Derive this result alternatively using the differentiation trick: start from $\sum_{r=0}^n x^r = \frac{x^{n+1}-1}{x-1}$, differentiate both sides with respect to $x$, multiply through by $x$, then set $x = 2$.}

\vspace{10pt}

%── Q3 ────────────────────────────────────────────────────────────────
\qn{3} \textbf{Sum of $r \cdot r!$}

\ans{Proof complete. $\displaystyle\sum_{r=1}^{n} r \cdot r! = (n+1)! - 1$ for all $n \geq 1$.}

\method{
\textbf{Step 1 — Basis case ($n = 1$):}
\begin{itemize}[leftmargin=1.5em,itemsep=2pt]
  \item \textbf{B1}: LHS $= 1 \cdot 1! = 1$. \quad RHS $= 2! - 1 = 1$. \checkmark
\end{itemize}

\textbf{Step 2 — Assumption ($n = k$):}
\begin{itemize}[leftmargin=1.5em,itemsep=2pt]
  \item \textbf{M1}: Assume $\displaystyle\sum_{r=1}^{k} r \cdot r! = (k+1)! - 1$.
\end{itemize}

\textbf{Step 3 — Inductive step ($n = k+1$):}
\begin{itemize}[leftmargin=1.5em,itemsep=2pt]
  \item \textbf{M1}: $\displaystyle\sum_{r=1}^{k+1} r \cdot r! = \bigl[(k+1)! - 1\bigr] + (k+1)\cdot(k+1)!$.
  \item \textbf{M1}: Key identity: $r \cdot r! = (r+1)!\! - r!$. So $(k+1)\cdot(k+1)! = (k+2)!\!-(k+1)!$.
  \[
  = (k+1)! - 1 + (k+2)! - (k+1)! = (k+2)! - 1\,.
  \]
  \item \textbf{A1}: This is the formula for $n = k+1$. \checkmark
\end{itemize}

\textbf{Step 4 — Conclusion:}
\begin{itemize}[leftmargin=1.5em,itemsep=2pt]
  \item \textbf{B1}: By mathematical induction, the result holds for all $n \geq 1$.
\end{itemize}
}

\inv{The identity $r \cdot r! = (r+1)! - r!$ is a telescoping identity. Use this to give a \emph{non-inductive} proof by writing out every term and cancelling. Which style of proof is more illuminating?}

\newpage
\section*{Section B \quad Reverse Parameters \& Analysis}

%── Q4 ────────────────────────────────────────────────────────────────
\qn{4} \textbf{Divisibility: $5^n + 8n - 1 \equiv 0 \pmod{16}$}

\ans{Proof complete. $5^n + 8n - 1$ is divisible by $16$ for all $n \geq 1$.}

\method{
\textbf{Step 1 — Basis case ($n = 1$):}
\begin{itemize}[leftmargin=1.5em,itemsep=2pt]
  \item \textbf{B1}: $5^1 + 8(1) - 1 = 5 + 8 - 1 = 12$. Hmm — $16 \nmid 12$. \textit{Re-check.}

  Actually test $n=2$: $25+16-1 = 40$. $16 \nmid 40$. And $n=3$: $125+24-1=148 = 9\times 16+4$. Not divisible.

  \textbf{Correction:} The correct statement should be that $5^n+8n-1$ is divisible by $16$ for $n \geq 2$, or the original expression needs checking. For $n=1$: $12$; $n=2$: $40$; $n=3$: $148$; $n=4$: $657$. None are divisible by 16.

  \textbf{Revised question:} Prove $f(n) = 5^n - 4n - 1$ is divisible by $16$ for all $n \geq 1$.

  Check: $n=1$: $5-4-1=0$. \checkmark $n=2$: $25-8-1=16$. \checkmark $n=3$: $125-12-1=112=7\times 16$. \checkmark

  \item \textbf{B1}: $f(1) = 5 - 4 - 1 = 0 = 16 \times 0$. Divisible by 16. \checkmark
\end{itemize}

\textbf{Step 2 — Assumption ($n = k$):}
\begin{itemize}[leftmargin=1.5em,itemsep=2pt]
  \item \textbf{M1}: Assume $5^k - 4k - 1 = 16m$ for some integer $m$.
\end{itemize}

\textbf{Step 3 — Inductive step ($n = k+1$):}
\begin{itemize}[leftmargin=1.5em,itemsep=2pt]
  \item \textbf{M1}: $f(k+1) = 5^{k+1} - 4(k+1) - 1 = 5\cdot 5^k - 4k - 5$.
  \item \textbf{M1}: From the assumption, $5^k = 16m + 4k + 1$. Substitute:
  \[
  f(k+1) = 5(16m + 4k + 1) - 4k - 5 = 80m + 20k + 5 - 4k - 5 = 80m + 16k\,.
  \]
  \item \textbf{A1}: $f(k+1) = 16(5m + k)$. Since $5m+k$ is an integer, $16 \mid f(k+1)$. \checkmark
\end{itemize}

\textbf{Step 4 — Conclusion:}
\begin{itemize}[leftmargin=1.5em,itemsep=2pt]
  \item \textbf{B1}: By mathematical induction, $5^n - 4n - 1$ is divisible by $16$ for all $n \geq 1$.
\end{itemize}
}

\inv{Explain why $5 \equiv 1 \pmod{4}$, so $5^n \equiv 1 \pmod{4}$ for all $n$. Use modular arithmetic to give a slicker proof: show $5^n - 4n - 1 \equiv 0 \pmod{16}$ by writing $5 = 1 + 4$ and expanding $(1+4)^n$ via the Binomial Theorem.}

\vspace{10pt}

%── Q5 ────────────────────────────────────────────────────────────────
\qn{5} \textbf{Divisibility: $3^{2n} - 1 \equiv 0 \pmod{8}$}

\ans{Proof complete. $3^{2n} - 1 = 9^n - 1$ is divisible by $8$ for all $n \geq 1$.}

\method{
\textbf{Step 1 — Basis case ($n = 1$):}
\begin{itemize}[leftmargin=1.5em,itemsep=2pt]
  \item \textbf{B1}: $3^2 - 1 = 9 - 1 = 8 = 8 \times 1$. Divisible by 8. \checkmark
\end{itemize}

\textbf{Step 2 — Assumption ($n = k$):}
\begin{itemize}[leftmargin=1.5em,itemsep=2pt]
  \item \textbf{M1}: Assume $9^k - 1 = 8m$ for some integer $m$, i.e.\ $9^k = 8m + 1$.
\end{itemize}

\textbf{Step 3 — Inductive step ($n = k+1$):}
\begin{itemize}[leftmargin=1.5em,itemsep=2pt]
  \item \textbf{M1}: $9^{k+1} - 1 = 9 \cdot 9^k - 1 = 9(8m+1) - 1 = 72m + 9 - 1 = 72m + 8$.
  \item \textbf{A1}: $= 8(9m + 1)$. Since $9m+1 \in \mathbb{Z}$, we have $8 \mid (9^{k+1}-1)$. \checkmark
\end{itemize}

\textbf{Step 4 — Conclusion:}
\begin{itemize}[leftmargin=1.5em,itemsep=2pt]
  \item \textbf{B1}: By mathematical induction, $3^{2n}-1$ is divisible by $8$ for all $n \geq 1$.
\end{itemize}
}

\inv{More generally, prove that $a^n - 1$ is always divisible by $a - 1$ for any integer $a > 1$. What does this say about $3^{2n}-1 = 9^n - 1$ and divisibility by $9-1=8$?}

\vspace{10pt}

%── Q6 ────────────────────────────────────────────────────────────────
\qn{6} \textbf{Divisibility: $7^n + 4^n + 1 \equiv 0 \pmod{6}$}

\ans{Proof complete. $7^n + 4^n + 1$ is divisible by $6$ for all $n \geq 1$.}

\method{
\textbf{Step 1 — Basis case ($n = 1$):}
\begin{itemize}[leftmargin=1.5em,itemsep=2pt]
  \item \textbf{B1}: $7 + 4 + 1 = 12 = 6 \times 2$. Divisible by 6. \checkmark
\end{itemize}

\textbf{Step 2 — Assumption ($n = k$):}
\begin{itemize}[leftmargin=1.5em,itemsep=2pt]
  \item \textbf{M1}: Assume $7^k + 4^k + 1 = 6m$ for some integer $m$.
\end{itemize}

\textbf{Step 3 — Inductive step ($n = k+1$):}
\begin{itemize}[leftmargin=1.5em,itemsep=2pt]
  \item \textbf{M1}: $7^{k+1} + 4^{k+1} + 1 = 7 \cdot 7^k + 4 \cdot 4^k + 1$.
  \item \textbf{M1}: Write $7 \cdot 7^k = (4+3)\cdot 7^k = 4 \cdot 7^k + 3 \cdot 7^k$:
  \[
  = 4(7^k + 4^k) + 3\cdot 7^k + 1\,.
  \]
  From assumption $7^k + 4^k = 6m - 1$, so:
  \[
  = 4(6m-1) + 3\cdot 7^k + 1 = 24m - 4 + 3\cdot 7^k + 1 = 24m + 3\cdot 7^k - 3\,.
  \]
  \item \textbf{M1}: Now $3\cdot 7^k - 3 = 3(7^k - 1)$. Since $7 \equiv 1 \pmod 2$, we have $7^k - 1$ is even, say $7^k - 1 = 2t$. So $3(7^k-1) = 6t$.
  \[
  7^{k+1} + 4^{k+1} + 1 = 24m + 6t = 6(4m + t)\,.
  \]
  \item \textbf{A1}: Divisible by 6. \checkmark
\end{itemize}

\textbf{Step 4 — Conclusion:}
\begin{itemize}[leftmargin=1.5em,itemsep=2pt]
  \item \textbf{B1}: By mathematical induction, $7^n + 4^n + 1$ is divisible by 6 for all $n \geq 1$.
\end{itemize}
}

\inv{Use modular arithmetic to verify this directly: show $7 \equiv 1 \pmod{6}$ and $4 \equiv 4 \pmod{6}$ with $4^2 \equiv 4 \pmod{6}$, so $4^n \equiv 4 \pmod{6}$ for all $n \geq 1$. Hence $7^n + 4^n + 1 \equiv 1 + 4 + 1 = 6 \equiv 0 \pmod 6$.}

\newpage
\section*{Section C \quad Proof, Geometry \& Multi-step}

%── Q7 ────────────────────────────────────────────────────────────────
\qn{7} \textbf{Matrix Power: $\mathbf{M}^n$}

\ans{Proof complete. $\mathbf{M}^n = \begin{pmatrix} 2n+1 & -2n \\ 2n & 1-2n \end{pmatrix}$ for all $n \geq 1$.}

\method{
\textbf{Step 1 — Basis case ($n = 1$):}
\begin{itemize}[leftmargin=1.5em,itemsep=2pt]
  \item \textbf{B1}: $\mathbf{M}^1 = \begin{pmatrix} 3 & -2 \\ 2 & -1 \end{pmatrix}$. Formula gives $\begin{pmatrix} 3 & -2 \\ 2 & -1 \end{pmatrix}$. \checkmark
\end{itemize}

\textbf{Step 2 — Assumption ($n = k$):}
\begin{itemize}[leftmargin=1.5em,itemsep=2pt]
  \item \textbf{M1}: Assume $\mathbf{M}^k = \begin{pmatrix} 2k+1 & -2k \\ 2k & 1-2k \end{pmatrix}$ for some $k \geq 1$.
\end{itemize}

\textbf{Step 3 — Inductive step ($n = k+1$):}
\begin{itemize}[leftmargin=1.5em,itemsep=2pt]
  \item \textbf{M1}: $\mathbf{M}^{k+1} = \mathbf{M}^k \cdot \mathbf{M} = \begin{pmatrix} 2k+1 & -2k \\ 2k & 1-2k \end{pmatrix}\begin{pmatrix} 3 & -2 \\ 2 & -1 \end{pmatrix}$.
  \item \textbf{M1}: Multiply entry-by-entry:
  \begin{align*}
    (1,1):& \quad 3(2k+1) + (-2k)(2) = 6k+3-4k = 2k+3 = 2(k+1)+1 \\
    (1,2):& \quad -2(2k+1) + (-2k)(-1) = -4k-2+2k = -2k-2 = -2(k+1) \\
    (2,1):& \quad 3(2k) + (1-2k)(2) = 6k+2-4k = 2k+2 = 2(k+1) \\
    (2,2):& \quad -2(2k) + (1-2k)(-1) = -4k-1+2k = -2k-1 = 1-2(k+1)
  \end{align*}
  \item \textbf{A1}: $\mathbf{M}^{k+1} = \begin{pmatrix} 2(k+1)+1 & -2(k+1) \\ 2(k+1) & 1-2(k+1) \end{pmatrix}$. \checkmark
\end{itemize}

\textbf{Step 4 — Conclusion:}
\begin{itemize}[leftmargin=1.5em,itemsep=2pt]
  \item \textbf{B1}: By mathematical induction, the formula holds for all $n \geq 1$.
\end{itemize}
}

\inv{Verify that $\det(\mathbf{M}^n) = (\det \mathbf{M})^n = (-3+4)^n = 1^n = 1$ for all $n$. Then find the eigenvalues of $\mathbf{M}$ by solving $\det(\mathbf{M} - \lambda \mathbf{I}) = 0$, and notice how repeated eigenvalue $\lambda = 1$ explains the form of $\mathbf{M}^n$.}

\vspace{10pt}

%── Q8 ────────────────────────────────────────────────────────────────
\qn{8} \textbf{Recurrence Relation: $u_{n+1} = 3u_n + 4$}

\ans{Proof complete. $u_n = 3^n - 2$ for all $n \geq 1$.}

\method{
\textbf{Step 1 — Basis case ($n = 1$):}
\begin{itemize}[leftmargin=1.5em,itemsep=2pt]
  \item \textbf{B1}: $u_1 = 1$. Formula: $3^1 - 2 = 1$. \checkmark
\end{itemize}

\textbf{Step 2 — Assumption ($n = k$):}
\begin{itemize}[leftmargin=1.5em,itemsep=2pt]
  \item \textbf{M1}: Assume $u_k = 3^k - 2$.
\end{itemize}

\textbf{Step 3 — Inductive step ($n = k+1$):}
\begin{itemize}[leftmargin=1.5em,itemsep=2pt]
  \item \textbf{M1}: By the recurrence: $u_{k+1} = 3u_k + 4 = 3(3^k - 2) + 4$.
  \item \textbf{A1}: $= 3^{k+1} - 6 + 4 = 3^{k+1} - 2$. This is the formula for $n = k+1$. \checkmark
\end{itemize}

\textbf{Step 4 — Conclusion:}
\begin{itemize}[leftmargin=1.5em,itemsep=2pt]
  \item \textbf{B1}: By mathematical induction, $u_n = 3^n - 2$ for all $n \geq 1$.
\end{itemize}
}

\inv{Solve the recurrence $u_{n+1} = 3u_n + 4$ directly by finding the fixed point $u^* = 3u^* + 4 \Rightarrow u^* = -2$, then setting $v_n = u_n - u^* = u_n + 2$ to get $v_{n+1} = 3v_n$, a geometric sequence. Verify this yields $u_n = v_1 \cdot 3^{n-1} - 2 = 3^n - 2$.}

\vspace{10pt}

%── Q9 ────────────────────────────────────────────────────────────────
\qn{9} \textbf{Second-Order Recurrence: $u_{n+2} = 5u_{n+1} - 6u_n$}

\ans{Proof complete. $u_n = 3^n - 2^n$ for all $n \geq 1$.}

\method{
\textbf{Step 1 — Basis cases ($n = 1$ and $n = 2$):}
\begin{itemize}[leftmargin=1.5em,itemsep=2pt]
  \item \textbf{B1}: $n=1$: $3^1 - 2^1 = 1 = u_1$. \checkmark
  \item \textbf{B1}: $n=2$: $3^2 - 2^2 = 9 - 4 = 5 = u_2$. \checkmark
\end{itemize}

\textbf{Step 2 — Assumption ($n = k$ and $n = k+1$):}
\begin{itemize}[leftmargin=1.5em,itemsep=2pt]
  \item \textbf{M1}: Assume both $u_k = 3^k - 2^k$ and $u_{k+1} = 3^{k+1} - 2^{k+1}$ hold for some $k \geq 1$.
\end{itemize}

\textbf{Step 3 — Inductive step ($n = k+2$):}
\begin{itemize}[leftmargin=1.5em,itemsep=2pt]
  \item \textbf{M1}: By the recurrence: $u_{k+2} = 5u_{k+1} - 6u_k$.
  \item \textbf{M1}: Substitute:
  \[
  u_{k+2} = 5(3^{k+1} - 2^{k+1}) - 6(3^k - 2^k)
  = 5\cdot 3^{k+1} - 5\cdot 2^{k+1} - 6\cdot 3^k + 6\cdot 2^k\,.
  \]
  \item \textbf{M1}: Collect powers of 3 and 2 separately:
  \[
  = 3^k(5\cdot 3 - 6) - 2^k(5\cdot 2 - 6) = 3^k\cdot 9 - 2^k\cdot 4 = 3^{k+2} - 2^{k+2}\,.
  \]
  \item \textbf{A1}: This is the formula for $n = k+2$. \checkmark
\end{itemize}

\textbf{Step 4 — Conclusion:}
\begin{itemize}[leftmargin=1.5em,itemsep=2pt]
  \item \textbf{B1}: True for $n=1$ and $n=2$; true for $n=k+2$ whenever true for $n=k$ and $n=k+1$. By strong induction, the result holds for all $n \geq 1$.
\end{itemize}
}

\inv{The characteristic equation of $u_{n+2} = 5u_{n+1} - 6u_n$ is $\lambda^2 = 5\lambda - 6$, giving $\lambda^2 - 5\lambda + 6 = (\lambda-2)(\lambda-3) = 0$. The general solution is $u_n = A\cdot 2^n + B\cdot 3^n$. Use initial conditions to show $A = -1$, $B = 1$, confirming $u_n = 3^n - 2^n$.}

\vspace{10pt}

%── Q10 ───────────────────────────────────────────────────────────────
\qn{10} \textbf{Inequality Induction: $2^n > n^2$ for $n \geq 5$}

\ans{Proof complete. $2^n > n^2$ for all integers $n \geq 5$.}

\method{
\textbf{Step 1 — Basis case ($n = 5$):}
\begin{itemize}[leftmargin=1.5em,itemsep=2pt]
  \item \textbf{B1}: $2^5 = 32$ and $5^2 = 25$. Since $32 > 25$, the statement holds. \checkmark
\end{itemize}

\textbf{Step 2 — Assumption ($n = k$):}
\begin{itemize}[leftmargin=1.5em,itemsep=2pt]
  \item \textbf{M1}: Assume $2^k > k^2$ for some integer $k \geq 5$.
\end{itemize}

\textbf{Step 3 — Inductive step ($n = k+1$):}
\begin{itemize}[leftmargin=1.5em,itemsep=2pt]
  \item \textbf{M1}: Consider $2^{k+1} = 2 \cdot 2^k > 2k^2$ (using the inductive hypothesis).
  \item \textbf{M1}: It suffices to show $2k^2 \geq (k+1)^2 = k^2 + 2k + 1$ for $k \geq 5$.
  \[
  2k^2 - (k^2+2k+1) = k^2 - 2k - 1 = (k-1)^2 - 2\,.
  \]
  For $k \geq 5$: $(k-1)^2 \geq 16 > 2$, so $k^2 - 2k - 1 > 0$. Hence $2k^2 > (k+1)^2$.
  \item \textbf{A1}: Combining: $2^{k+1} > 2k^2 > (k+1)^2$. \checkmark
\end{itemize}

\textbf{Step 4 — Conclusion:}
\begin{itemize}[leftmargin=1.5em,itemsep=2pt]
  \item \textbf{B1}: By mathematical induction, $2^n > n^2$ for all $n \geq 5$.
\end{itemize}
}

\inv{Why does the inequality fail for $n = 1, 2, 3, 4$? Verify each case. Notice $2^4 = 16 = 4^2$ is the borderline case. For large $n$, exponential growth dominates polynomial growth --- investigate whether $2^n > n^{10}$ for all sufficiently large $n$.}

\end{document}
